% Generated draft: requires mathematical and source review.
\documentclass[11pt]{article}
\usepackage[margin=25mm]{geometry}
\usepackage{fontspec}
\setmainfont{Latin Modern Roman}
\usepackage{amsmath,amssymb,mathtools,unicode-math}
\setmathfont{Latin Modern Math}
\usepackage{enumitem,needspace,xurl,hyperref}
\setlength{\parindent}{0pt}
\setlength{\parskip}{5pt}
\setlength{\emergencystretch}{3em}
\newcommand{\N}{\mathbb{N}}
\newcommand{\Z}{\mathbb{Z}}
\newcommand{\Q}{\mathbb{Q}}
\newcommand{\R}{\mathbb{R}}
\newcommand{\C}{\mathbb{C}}
\newcommand{\F}{\mathbb{F}}
\newcommand{\A}{\mathbb{A}}
\newcommand{\PP}{\mathbb{P}}
\newcommand{\E}{\mathbb{E}}
\newcommand{\eps}{\varepsilon}
\newcommand{\dd}{\,\mathrm{d}}
\begin{document}
\section*{UnsolvedMath ID: 30004346}
\textbf{Source problem code:} catalogue entry
\subsection{Definitions and mathematical statement}
The source uses the terms \emph{rigid meromorphic cocycle} \(J\), \emph{real-multiplication point} \(\tau\), and its value \(J[\tau]\), but does not specify their underlying arithmetic and analytic framework. Write \(H_\tau\) for the narrow ring class field attached to \(\tau\), as designated by the source. The associated real quadratic field, order, and field embeddings are not specified. For subfields \(K,L\) of a common ambient field, their compositum \(KL\) is the smallest subfield containing both.

The supplied conjectural assertion is: for every rigid meromorphic cocycle \(J\), there exists a finite extension \(H_J/\Q\) such that
\[
  J[\tau]\in H_JH_\tau
  \qquad\text{for every real-multiplication point }\tau.
\]
The extension \(H_J\) may depend on \(J\), but must work simultaneously for every \(\tau\). The assertion requires a common ambient field in which the values and the two number fields can be compared.

This is a qualified transcription rather than a complete intrinsic formulation: the source does not define the cocycle space, the domain of real-multiplication points, the evaluation and its normalization, or the order determining \(H_\tau\). In particular, it does not explain whether evaluations are defined and finite at every point quantified over, or how poles and other undefined evaluations are handled. These omissions require mathematical review; no restriction of the stated universal quantifier is imposed here.
\subsection{Short English statement}
The conjecture asserts that each rigid meromorphic cocycle has a single finite number field extension which, when composed with the narrow ring class field of any real-multiplication point, contains its value there. The supplied record leaves the precise cocycle and evaluation framework unspecified.
\subsection{Sources}
\begin{itemize}
\item Local UnsolvedMath registry, \texttt{data/unsolvedmath/problems.json}, permanent ID \(30004346\), source code \texttt{catalogue entry}, title: Algebraicity of Rigid Cocycle Values at Real Multiplication Points. The supplied record gives proposed year 2019 and names no individual proposer.
\item Supplied report citation: \emph{Modular Forms (2020)}, Oberwolfach Reports, pp.\ 3529--3606, DOI \texttt{10.4171/owr/2019/57}. \url{https://doi.org/10.4171/owr/2019/57}. The background attributes the extraction to workshop open-problem discussions documented in Oberwolfach Reports.
\item Darmon, rigid meromorphic cocycles: \url{https://www.math.mcgill.ca/darmon/pub/Articles/Expository/15.Rigid-Meromorphic-Cocycles/paper.pdf}.
\item Darmon--Pozzi--Vonk, Dedekind--Rademacher cocycle: \url{https://doi.org/10.4171/JEMS/1344}.
\item Literature assessment dated \texttt{2026-08-22}: supplied dated curation, not independently verified. It labels the full assertion for arbitrary rigid meromorphic cocycles and all real-multiplication points as open, and reports algebraicity results for special Dedekind--Rademacher values rather than the full extracted assertion. The background describes this assessment as literature triage derived from the supplied OWR curation and requests independent verification of time-sensitive status and cited claims.
\end{itemize}
\end{document}
